\begin{abstract}
For the cusp-normalized $(3,4,\infty)$ family on $S^6$, the particular period is $\frac{6i}{\pi}\operatorname{Log}\frac{\vartheta_1(\mu,\tau)}{2\sqrt3\,\eta_D(\tau)}$; its real determinant function has maximum $-(3/\pi)\log6$, attained only at the multiplicity-four fibre. The known punctured-disc parameter $w=e^{-\pi ic}$ therefore has exact radius $216$, and the Albanese nomes are $iw/216$ and $-w^2/48^3$. The period construction also gives the stated all-data invariants, Hodge and tangent-twist calculations, real structures, index-$12$ monodromy group, and descriptions of the special fibres. At the outer circle the family has no extension of the specified kind; Engel's proposed central space has virtual Euler number $0$ instead of $2$. Theorems are assembled from proofs by GPT-6 Astra (ChatGPT Pro mode), with the bounded readings and computations of Claude Fable 5.2 (efforts high and max) and Claude Opus 5.5 (effort max) specified per result; no human has checked them. Announcements retain repository proofs, including the zero-end non-filling, trivial smooth monodromy and component claims. The construction is credited to Claude (effort not recorded) and Alp\"oge and to Engel; the intrinsic pencil, standard-data $2\mathbb Z$ theorem, tangent cohomology, Kuranishi germ and full automorphisms are credited to Du, Guo, Kusakabe and Xie, arXiv:2610.00250v1.
\end{abstract}

\section*{Result}
\subsection*{Statement}
\noindent The rows below locate the main groups of results. Each individual statement, its status, its proof and its kind have a separate row in Appendix~\ref{app:audits}; its run identifier is retained in its heading. Here T means a theorem with a printed proof, K a known theorem re-derived in Appendix~\ref{app:known}, and A an announcement with proof in the repository.
\begin{center}\small
\begin{tabularx}{\linewidth}{@{}lYll@{}}
\toprule Identifier & Statement & Status & Proof / kind\\\midrule
FAM B.1--B.3 & Dictionary $u=w^2$; all discrete data; $2kl\mathbb Z$ in the stated model & T & \S\S1--2 / theorem\\
FAM C; ARI C & Closed periods, division values, maximum and radius $216$ & T & \S\S3--4 / theorem, certificate\\
INV A--D; CDP C5b & Hodge, Fr\"olicher, Bott--Chern, Aeppli and tangent twists & T & \S\S5--6 / theorem\\
FAM B.4, D, F & Real structures, monodromy, special fibres and subvarieties & T & \S\S7--9 / theorem\\
PUSH D.1, A.3 & Outer non-extension; conjectured centre has the wrong Euler number & T & \S10 / counterexample\\
DGKX v1 & Intrinsic pencil, $2\mathbb Z$, deformations, $\Aut=\mathbb C^*$ & K & App. B / new proof\\
Remaining selected rows & Exact announcement statements and their mechanisms & A & \S\S2--7,10--11 / theorem\\
\bottomrule
\end{tabularx}
\end{center}
For example, $c=0$ gives $w=1$ and Albanese nome $i/216$. The component and zero-end assertions remain announcements; the dictionary for higher multiplicities is not included in the theorem for the stated model.

\subsection*{Status}
The same model, GPT-6 Astra (ChatGPT Pro mode), produced the six mathematical returns and assembled this paper. Claude Fable 5.2 (effort high) reviewed FAM and read the explicitly listed parts of the intakes; Claude Fable 5.2 (effort max) supplied the expert readings and derivations, and Claude Opus 5.5 (effort max) supplied the main-session checks recorded there. The historical program checks include the exact recomputations, the full interval re-execution and the Arb certification. This assembly adds source comparisons and opened-citation checks, not a second independent mathematical review. No human has checked any of it.
\begin{center}\small
\begin{tabularx}{\linewidth}{@{}lY@{}}
\toprule Kind & Theorem, computation, correction, counterexample and new proof, assigned per result.\\
Claim & T: proved as in the supplied text; K: known and re-derived; A: announced.\\
Checks & One row per recorded check in Appendix~\ref{app:verification}; scope and limitations are explicit.\\
Human understanding & none.\\
Formalization & none.\\
Open points & Four question groups in \S\ref{sec:open}; source discrepancies in Limits and open points.\\\bottomrule
\end{tabularx}
\end{center}

\subsection*{Known and new}
The section-by-section credit table in Appendix~\ref{app:credits} starts from the expert's supplied table and incorporates the three-source audit. We credit \cite[Main Theorem; Sections 3--9]{AC26} for the construction, canonical bundle and the stated cohomological, fibre and automorphism results, \cite[Sections 2--8]{E26} for the geometric construction, and \cite[Theorems 1.3, 1.4, 7.1 and 7.2]{DGKXv1} for the known deformation and period-identification results. In particular, the N\'eron--Severi equation of \cite[Lemma 9.2]{AC26} and the elliptic-logarithm reading of \cite[Remark 3.12]{AC26} are credited there, although their source-run proofs are retained. The source audit used the supplied full TeX and transcription and the primary references opened on 16 November 2026; negative searches are bounded searches, not assertions of absence from all literature.

The constructed $X_c$ has algebraic dimension $1$ and $c_3(X_c)=2$. This contradicts Corollary~2.3 of the corrigendum of Campana, Demailly and Peternell \cite[Corollary 2.3]{CDP20}, which asserts algebraic dimension $0$ for a complex structure on $S^6$. The counterexample is due to \cite[Main Theorem; Section 10]{AC26}, whose Section~10 locates two failing steps for $X$. Where the argument fails in general, and what is true in its place, is the subject of a separate note.

\subsection*{Proof map}
\textbf{Idea.} The marked period lattice carries the affine deck data, the theta norm and the root-base direct images. Its norm determines the admissible disc; its two finite fibres retain the Albanese nomes; its cusp supplies the non-normal fibre used in the cohomology calculations. The exact period identification is a known input.
\begin{enumerate}[leftmargin=*,label=\arabic*.]
\item Fix the marked construction and dictionary; use the norm obstruction for affine data (\S\S1--2, Appendix B).
\item Compute the theta expression and division values; evaluate the maximum by the certificates and ladder (\S\S3--4).
\item Calculate direct images on the root base and pass to the requested cohomology (\S\S5--6).
\item Read real forms, monodromy and fibres from the same lattice, and test the boundary and central model (\S\S7--10).
\end{enumerate}
\textbf{Hard step.} The exact maximum uses the two interval certifications as proof of record. The supplied review records a full re-execution of the $19\,701$ norm boxes and a derivation of the local estimate; the ladder's displayed inequalities were checked, but ARI C.8.1--C.8.3 were not read line by line in that intake. The certificate files were not supplied to this assembly.

\textbf{Reading routes.} The results table and credit table give the claims and their scope. A reader using the known period-identification theorem can read \S\S1--4 and then the relevant application. A full check follows the dependency rows in Appendix~\ref{app:dependencies}, including the printed re-derivations in Appendices A--B. The principal external inputs, with exact locators and S-checks, are the cited construction, the known DGKX theorems, the Mordell--Weil height formula, proper base change and relative duality, the structure theorem for double complexes, and the extension results used in \S10.
\clearpage
